\documentclass[11pt]{article}
\usepackage[margin=0.85in]{geometry}
\usepackage{booktabs}
\usepackage{amsmath}
\usepackage{parskip}
\usepackage{xcolor}
\usepackage{titlesec}

\definecolor{e3blue}{HTML}{2A78D6}
\titleformat{\section}{\large\bfseries\color{e3blue}}{\thesection}{0.5em}{}

\title{\textbf{Q2: 2023 Backcast Revenue Calculation}\\
\large GCV Solar + Storage Valuation --- NYISO N.Y.C. Zone}
\author{E3 Applicant Exercise}
\date{2026}

\begin{document}
\maketitle

\section*{Inputs and Notation}

All revenue is calculated from hourly 2023 NYISO N.Y.C. zone real-time prices.
For hour $t$, let $p_t$ be the real-time energy price (\$/MWh) and $s_t$ be the
normalized solar production shape. Solar output is assumed to be available for
the full hour, so MW and MWh are numerically equivalent for each hourly interval.

\begin{table}[h!]
\centering
\begin{tabular}{ll}
\toprule
\textbf{Parameter} & \textbf{Value} \\
\midrule
Solar size & 300 MW-DC \\
Solar inverter limit & 250 MW-AC \\
Storage power capacity & 50 MW \\
Storage energy capacity & 200 MWh \\
Storage duration & $200\text{ MWh} / 50\text{ MW} = 4$ hours \\
Round-trip efficiency & 100\% per provided input \\
Dispatch simplification & One full charge/discharge cycle per day with perfect real-time foresight \\
\bottomrule
\end{tabular}
\end{table}

\section*{Q2a. Standalone Solar Revenue}

Hourly solar generation is calculated from the normalized production shape and
capped at the inverter limit:

\[
G_t = \min(s_t \times 300\text{ MW},\ 250\text{ MW})
\]

Annual backcast solar revenue is the sum of hourly delivered energy times the
real-time price:

\[
R_{\text{solar}} = \sum_{t=1}^{8760} G_t \times p_t
\]

Using the 2023 hourly data:

\[
\sum_t G_t = 551{,}454\text{ MWh}
\]

\[
R_{\text{solar}} = \$18{,}910{,}425
\]

\[
\text{Solar capture price}
= \frac{\$18{,}910{,}425}{551{,}454\text{ MWh}}
= \$34.29/\text{MWh}
\]

The inverter cap clips 372 MWh, equal to 0.07\% of gross solar generation, so
clipping has a negligible effect on standalone solar revenue.

\section*{Q2b. Standalone Storage Revenue}

Storage is modeled as daily energy arbitrage. For each day $d$, the battery
charges during the four lowest-priced hours and discharges during the four
highest-priced hours:

\[
R_{\text{storage},d}
= 50\text{ MW} \times
\left(
\sum_{h \in H^{\text{high}}_d} p_h
-
\sum_{h \in H^{\text{low}}_d} p_h
\right)
\]

where $H^{\text{high}}_d$ is the set of the day's four highest-price hours and
$H^{\text{low}}_d$ is the set of the day's four lowest-price hours.

Annual standalone storage revenue is:

\[
R_{\text{storage}}
= \sum_{d=1}^{365} R_{\text{storage},d}
= \$2{,}473{,}753
\]

Normalized by installed capacity:

\[
\frac{\$2{,}473{,}753}{50{,}000\text{ kW}}
= \$49.48/\text{kW-year}
\]

\[
\frac{\$2{,}473{,}753}{200{,}000\text{ kWh}}
= \$12.37/\text{kWh-year}
\]

\section*{Q2c. Combined Solar + Storage Revenue}

The combined case keeps the same solar export calculation as standalone solar.
Storage first charges from otherwise-clipped solar energy at zero opportunity
cost, then charges any remaining daily energy requirement from the grid during
low-price hours.

For each hour, clipped solar is:

\[
C_t = \max(s_t \times 300\text{ MW} - 250\text{ MW},\ 0)
\]

Daily storage charging from clipped solar is limited by storage power and daily
energy capacity:

\[
E^{\text{clip}}_d
= \min\left(
200\text{ MWh},\
\sum_{t \in d} \min(C_t,\ 50\text{ MW})
\right)
\]

The remaining daily storage charge is bought from the grid:

\[
E^{\text{grid}}_d = 200\text{ MWh} - E^{\text{clip}}_d
\]

Storage then discharges during the day's four highest-price hours, as in the
standalone storage case. Across 2023, the model uses:

\[
\sum_d E^{\text{clip}}_d = 372\text{ MWh}
\]

\[
\sum_d E^{\text{grid}}_d = 72{,}628\text{ MWh}
\]

The resulting combined revenues are:

\[
R_{\text{combined solar}} = \$18{,}910{,}425
\]

\[
R_{\text{combined storage}} = \$2{,}480{,}611
\]

\[
R_{\text{combined total}}
= \$18{,}910{,}425 + \$2{,}480{,}611
= \$21{,}391{,}036
\]

The interaction benefit versus building standalone solar and standalone storage
separately is:

\[
R_{\text{combined total}} -
\left(R_{\text{solar}} + R_{\text{storage}}\right)
= \$21{,}391{,}036 - (\$18{,}910{,}425 + \$2{,}473{,}753)
= \$6{,}857
\]

\[
\frac{\$6{,}857}{\$21{,}384{,}178}
= 0.032\%
\]

\section*{Summary}

\begin{table}[h!]
\centering
\begin{tabular}{lp{1.75in}p{3.0in}}
\toprule
\textbf{Configuration} & \textbf{2023 Backcast Revenue} & \textbf{Key Method} \\
\midrule
Standalone solar & \$18.910M & Hourly solar output $\times$ hourly RT price \\
Standalone storage & \$2.474M & Daily four-low / four-high arbitrage \\
Combined solar + storage & \$21.391M & Solar revenue plus storage with clipped solar charging first \\
\bottomrule
\end{tabular}
\end{table}

The combined case produces only a small energy-arbitrage uplift over the
standalone sum because the solar system clips very little energy in 2023.

\end{document}
